\section{Conditional Statements and Decision Trees}

A program that only runs from top to bottom can describe a fixed procedure, but it cannot react to different runtime states. Conditional statements introduce decision points into the execution graph. At each decision point, Python evaluates a condition and selects one route while skipping the others.

This section studies how those branching routes are written, nested, and evaluated. The visible syntax of \texttt{if}, \texttt{elif}, and \texttt{else} defines the shape of the decision tree, while Python's truth-value rules decide which edge is followed at runtime.

The central idea is that a condition is not merely a decorative test placed before some code. It is a routing mechanism. It decides whether a block becomes reachable, whether a nested decision is even inspected, and whether execution continues through one path or another.